% operator-stability-v1 — publication template (writeup-aligned)
% Slots filled by scripts/generate_operator_paper.py
\documentclass[12pt]{article}
\usepackage[letterpaper,top=2cm,bottom=2cm,left=3cm,right=3cm]{geometry}
\usepackage[T1]{fontenc}
\usepackage[utf8]{inputenc}
\usepackage{lmodern}
\usepackage{microtype}
\usepackage{amsmath,amssymb,amsthm,mathtools}
\usepackage{booktabs}
\usepackage{xcolor}
\usepackage{enumitem}
\usepackage{xurl}
\usepackage[colorlinks=true,linkcolor=blue,urlcolor=blue,citecolor=blue]{hyperref}

\Urlmuskip=0mu plus 1mu
\urlstyle{tt}
\emergencystretch=2em

\setlength{\parskip}{0.4em}
\setlength{\parindent}{1.2em}
\setlist[itemize]{leftmargin=1.4em,itemsep=0.25em,topsep=0.35em}
\setlist[enumerate]{leftmargin=1.4em,itemsep=0.25em,topsep=0.35em}

\theoremstyle{plain}
\newtheorem{theorem}{Theorem}
\newtheorem{lemma}[theorem]{Lemma}
\theoremstyle{definition}
\newtheorem{definition}[theorem]{Definition}
\newtheorem{example}[theorem]{Example}
\theoremstyle{remark}
\newtheorem{remark}[theorem]{Remark}

% Breakable monospace for Lean paths / identifiers (xurl \path)
\newcommand{\leanpath}[1]{%
  {\raggedright\urlstyle{tt}\path{#1}}}


\definecolor{fundbox}{RGB}{245,245,248}
\definecolor{fundrule}{RGB}{180,180,190}

\title{Partial sorting Ranking Preservation under Bounded Score Perturbations}
\author{Research Architecture Runtime\\[0.25em]
{\normalsize\texttt{operator/partial-sorting}}}
\date{2026-07-29}

\begin{document}
\maketitle

\begin{abstract}
\leanpath{partial-sorting} is certified by pairwise ranking preservation when all gaps exceed $2\varepsilon$. Lean \texttt{LEAN\_FULL} via \texttt{OrderStat.Ranking}.
\end{abstract}

\noindent
\begin{center}
\fcolorbox{fundrule}{fundbox}{%
  \begin{minipage}{0.92\linewidth}
  \centering
  \textbf{This theorem is:} \texttt{Reduction}
  \end{minipage}}
\end{center}
\vspace{0.6em}

\section{Problem}
{\raggedright\sloppy
The \leanpath{partial-sorting} operator depends on the pairwise order of scores $s\in\mathbb{R}^n$. The certified object is the full pairwise ranking (and thereby any ranking-determined selection such as top-$k$ index sets).
\par}

\section{Stability notion}
Perturbations: $\|\delta\|_\infty\le\varepsilon$. Stability: if $\mathrm{AllGapsExceed}(s,2\varepsilon)$, then $s_i<s_j\Leftrightarrow (s+\delta)_i<(s+\delta)_j$ for all $i,j$. Sharpness: a gap $\le 2\varepsilon$ admits a collision adversary.

\section{Definitions}
\begin{definition}[All-gaps]
$|s_p-s_q|>g$ for $p\neq q$.\end{definition}
\noindent\textbf{Assumptions.} $n\ge 2$, $\varepsilon\ge 0$, and (for invariance) all pairwise gaps $>2\varepsilon$.

\section{Theorem}
\begin{theorem}[Ranking invariance]\label{thm:inv}
If $\mathrm{AllGapsExceed}(s,2\varepsilon)$, then for every $\|\delta\|_\infty\le\varepsilon$ and all $i,j$, $s_i<s_j\Leftrightarrow s_i+\delta_i<s_j+\delta_j$.
\end{theorem}
\begin{theorem}[Sharpness]\label{thm:sharp}
If some distinct pair has $|s_i-s_j|\le 2\varepsilon$, an admissible $\delta$ forces a value collision.
\end{theorem}

\section{Intuition}
Relative $\ell_\infty$ moves cost at most $2\varepsilon$. Gaps larger than $2\varepsilon$ cannot reverse; smaller gaps can be closed by a midpoint push.

\section{Examples}
\begin{example}
Scores $(0,3,8)$, $\varepsilon=1$: all gaps $>2$, so the ranking is stable. If scores are $(0,1.5)$ and $\varepsilon=1$, the gap $1.5\le 2$ admits a midpoint tie.
\end{example}

\section{Proof}
Fix $\|\delta\|_\infty\le\varepsilon$. If $i=j$ both sides of the claimed equivalence are false for strict $<$. If $i\neq j$ and $|s_i-s_j|>2\varepsilon$, the estimate $(s_i+\delta_i)-(s_j+\delta_j)\ge (s_i-s_j)-2\varepsilon$ (and its swap) yields $s_i<s_j\Leftrightarrow s_i+\delta_i<s_j+\delta_j$ (Theorem~\ref{thm:inv}).

For sharpness, take $i\neq j$ with $|s_i-s_j|\le 2\varepsilon$ and the midpoint $\delta_i=-(s_i-s_j)/2$, $\delta_j=(s_i-s_j)/2$ (elsewhere $0$). Then $\|\delta\|_\infty\le\varepsilon$ and $s_i+\delta_i=s_j+\delta_j$ (Theorem~\ref{thm:sharp}).

\section{Formal statement}
{\raggedright\sloppy
\begin{description}[style=nextline,leftmargin=1.1em,font=\normalfont\bfseries]
\item[Lean module] \leanpath{Research.Operators.PartialSorting.Preservation}
\item[Source file] \leanpath{lean/Research/Operators/PartialSorting/Preservation.lean}
\item[Theorems]\leavevmode\par\vspace{0.15em}\begin{itemize}[leftmargin=1.2em]
  \item \leanpath{partial\_sorting\_margin\_invariance}
  \item \leanpath{partial\_sorting\_margin\_sharpness}
\end{itemize}
\item[Note] Reduction of OrderStat.Ranking.
\item[Certificate] \leanpath{lean/certificates/partial-sorting/partial-sorting-margin}
\item[Status] \texttt{LEAN\_FULL} on Mathlib $\mathbb{R}$ (\texttt{REAL\_MATHLIB})
\end{description}
\par}

\section{Proof dependencies}
{\raggedright\sloppy
\begin{itemize}
\item \leanpath{Research.Operators.OrderStat.Ranking} / shared pairwise lemma from \leanpath{KthMargin}.
\item Midpoint collision adversary.
\item No other operator theorems required.
\end{itemize}
\par}

\section{Consequences}
{\raggedright\sloppy
Ranking stability implies stability of any functional of the pairwise order used by \leanpath{partial-sorting}. Related operators: \leanpath{top-k}, \leanpath{sorting}, \leanpath{rank}.
\par}

\end{document}
